% \section{Complexity analysis}
% \label{Section: Convergence analysis}

% We now investigate the complexity analysis for DualHash-StoM (\Cref{alg:our_algorithm_sgdm}) and DualHash-StoRM (\Cref{alg:our_algorithm_storm}). Our objective is to find an approximate stochastic critical point of the Lagrangian function \eqref{Eq: simplified_lagrangian}, i.e., a point satisfying 
% \begin{equation}\label{equ:esp}
% \mathbb{E} [\text{dist}^2(\bm{0}, \partial \calL (\bx^*, \bm{B}^*, \bm{\Lambda}^* ))] \leq \varepsilon^2
% \end{equation}
% for a given \(\varepsilon > 0\). We assume in this paper the set \(
% \text{crit} \calL := \{(\bx^*, \bm{B}^*, \bm{\Lambda}^* ) \in \mathcal{X} \times \mathcal{B} \times \mathcal{Y^*}: \bm{0} \in \partial \mathcal{L}(\bx^*, \bm{B}^*, \bm{\Lambda}^*)\}
% \) is nonempty. This focus is standard in non-convex stochastic optimization analysis, as finding a global or even local minimizer for non-convex optimization problems is generally NP-hard. In particular, we can prove that a critical point of the Lagrangian function \eqref{Eq: simplified_lagrangian} is a first-order necessary optimal condition of problem \eqref{Prop: block-wise regularized problem} as follows. 

% \begin{lemma}
% \label{Lemma: Optimality condition}
%     Suppose that \(\min\limits_{\bx, \bB} \calL(\bx, \bB, \bLmbd) > -\infty\)  and \((\bx^*, \bB^*, \bLmbd^*)\) are the optimal solutions of the problem \eqref{Prop: block-wise regularized problem} and the problem \eqref{Prop: the dual problem}, respectively. Then, it holds that \((\bx^*, \bB^*, \bLmbd^*) \in \text{crit} \calL\), i.e., \(\bm{0} \in \partial \mathcal{L}(\bx^*, \bB^*, \bLmbd^*)\).
% \end{lemma} 
% \setcounter{proof}{0}
% \begin{proof}
%     According to \cite[Theorem 2.158]{temporary-citekey-11039},
%     $(\bx^{*}, \bm{B}^{*}, \bm{\Lambda}^{*})$ are optimal solutions to the
%     primal problem \eqref{Prop: block-wise regularized problem} and the dual problem
%     \eqref{Prop: the dual problem} if and only if:
%     \begin{align}
%         \label{Eq: the sufficient and necessary condition}
%         \begin{cases}(\bx^{*}, \bm{B}^{*}) \in \argmin\limits_{\bx \in \calX, \bm{B} \in \mathcal{B}}\mathcal{L}(\bx^{*}, \bm{B}^{*}, \bm{\Lambda}^{*}), \\ 0 = h(\bm{B}^{*}) + h^{*}(\bm{\Lambda}^{*}) - \langle \bm{\Lambda}^{*}, \bm{B}^{*}\rangle.\end{cases}
%     \end{align}

%     The first condition in \eqref{Eq: the sufficient and necessary condition} implies
%     that $\bm{0}$ belongs to the subdifferential of $\mathcal{L}$ with respect to
%     $(\bx, \bm{B})$ at $(\bx^{*}, \bm{B}^{*}, \bm{\Lambda}^{*})$. Since $\calL$
%     is continuously differentiable in $(\bx, \bm{B})$, we have
%     \begin{align}
%         \label{Eq: the sufficient and necessary condition for primal}\begin{cases}\bm{0}&= \nabla f(\bx^{*}) + \nabla_{\bx}G(\bx^{*}, \bm{B}^{*}),\\ \bm{0}&= \nabla_{\bB}G(\bx^{*}, \bm{B}^{*}) + \bm{\Lambda}^{*}.\end{cases}
%     \end{align}
%     For the second condition in \eqref{Eq: the sufficient and necessary condition},
%     by the definition of the conjugate function, if this equality holds, then $\bm
%     {\Lambda}^{*}$ is the optimal solution of $\sup \{ \left<\bb^{*}, \blmbd\right
%     > - h^{*}(\blmbd) \}$, i.e.,
%     \begin{align}
%         \label{Eq: the sufficient and necessary condition for dual}\bm{B}^{*}\in \partial h^{*}(\bm{\Lambda}^{*}).
%     \end{align}
%     Combining \eqref{Eq: the sufficient and necessary condition for primal} and \eqref{Eq: the sufficient and necessary condition for dual},
%     we obtain $\bm{0}\in \partial \mathcal{L}(\bx^{*}, \bm{B}^{*}, \bm{\bLmbd}^{*}
%     )$, completing the proof. 
%     %, which can be written explicitly as:
%     % \begin{align}
%     %     \label{Eq: necessary_condition_dual}
%     %     \begin{cases}
%     %         \bm{0} &= \nabla f(\bx^*) + \nabla_{\bx} G(\bx^*, \bm{B}^*),\\
%     %         \bm{0} &= \nabla_{\bB} G(\bx^*, \bm{B}^*) + \bm{\Lambda}^*,\\
%     %         \bm{0} &\in \bm{B}^* - \partial h^*(\bm{\Lambda}^*).
%     %     \end{cases}
%     % \end{align}
%     % This system of equations characterizes the critical points of $\mathcal{L}$, thus completing the proof.
% \end{proof}

% We then make the following assumptions generally used in stochastic optimization \cite{Ghadimi2013MinibatchSA,huangMiniBatchStochasticADMMs2019,xu2023momentum}.
% \begin{assumption} 
% \label{Ass: main assumptions}
% \begin{itemize}[leftmargin=2em]
%     \item[{(i)}] We assume that the sequences \( \{(\bx^k, \bm{B}^k, \bm{\Lambda}^k) \}_{k \in \NN} \) generated by Algorithms \ref{alg:our_algorithm_sgdm} and \ref{alg:our_algorithm_storm} are bounded. Moreover, the objective of  \eqref{Prop: block-wise regularized problem} is lower bounded by \(J^*\).
%     \item[{(ii)}] We assume that \(F(\bx, \bB)\) is $L_F$-smooth with respect to $(\bx, \bB)$, where $L_F = \mathcal{O}\left(1\right)$.
%     \item[{(iii)}] Suppose index $i \sim \text{Uniform}\{1, \ldots, n\}$.  For any \(\bB \in \calB\), there exists a constant $\sigma^2 \geq 0$ such that $\mathbb{E}_{i}[\nabla_{\bx} F_i(\bx, \bb_i)] = \nabla_{\bx} F(\bx, \bB)$ and $\mathbb{E}_{i}[\|\nabla_{\bx} F_i(\bx, \bb_i) - \nabla_{\bx} F(\bx, \bB)\|^2] \leq \sigma^2$. 
% \end{itemize}
% \end{assumption}
% \begin{remark}
% \begin{enumerate}[label=(\roman*), leftmargin=2em]
%     \item The boundedness in \Cref{Ass: main assumptions} (i) is maintained in practice for each variable:  $\bm{\Lambda}^k$ remains bounded due to its proximal update rule \eqref{Eq: the update of lambda} for $\bm{\Lambda}^k$ with the truncation property of $\text{prox}_{h^*}$ in \eqref{Eq: the proximal mapping of conjugate function}; The $\bB^k$ naturally converges toward $\pm 1$ through W-type regularization; and network parameters $\bx^k$ are stabilized via weight decay and pretrained initialization, ensuring objective coercivity.
%     \item The smoothness constant \(L_F\) remains \(\mathcal{O}(1)\), as the penalty formulation is designed to cancel the potential \(\mathcal{O}(\sqrt{n})\) scaling of individual components. A rigorous analysis is provided in Appendix C.
% \end{enumerate}
% \end{remark}
% Our analysis aims to bound \( \EE [\| \partial \calL\|^2]\). However, due to the existence of dual variables \(\bLmbd\), the Lagrangian function \eqref{Prop: the dual problem} lacks the descent property. This makes it difficult to apply standard approaches from prior work \cite{tran2022hybrid,wang2019spiderboost,xu2023momentum,zhang2022stochastic}, which often leverage the descent of their objective functions or corresponding potential functions. Instead, we shall construct an appropriate Lyapunov function. To enable this construction, we first analyze how to control the dual variable dynamics by establishing bounds on \(\|\bm{\Lambda}^{k+1} - \bm{\Lambda}^k \|^2\) with the primal variables residuals.
% \begin{lemma}
% \label{Lemma: dual boundedness}
% Under \Cref{Ass: main assumptions}, for \(k \geq 1 \), it holds that
%     \begin{align}
%        \nonumber
%         &\|\bm{\Lambda}^{k+1}- \bm{\Lambda}^k \|^2 \\
%         \leq&\, 3 \tau^{-2} \|\bB^{k+2} - \bB^{k+1}\|^2 + 3(\tau^{-1} + L_F)^2  \|\bB^{k+1} - \bB^{k}\|^2   \notag \\
%         &+ 3L_F^2  \|\bx^{k+2} - \bx^{k+1}\|^2.
%          \label{Eq: dual boundedness}
%     \end{align}
% \end{lemma}
% \begin{proof}
%     Using the update of $\bm{B}^{k}$ in \eqref{Eq: the update of B} twice yields
%     that
% \begin{align*}
%      & \|\bm{\Lambda}^{k+1}- \bm{\Lambda}^{k}\| \notag    \\
%         = &\,\left\| \tau^{-1}(\bm{B}^{k+1}- \bm{B}^{k+2}) + \tau^{-1}(\bm{B}^{k+1}- \bm{B}^{k}) \right. \\ 
%         & \left. \quad + \nabla_{\bB}F(\bx^{k+1}, \bm{B}^{k}) - \nabla_{\bB}F(\bx^{k+2}, \bm{B}^{k+1}) \right\| \\
%        \leq & \,\tau^{-1}\|(\bm{B}^{k+1}- \bm{B}^{k+2})\| + (\tau^{-1}+ L_{F}) \|\bm{B}^{k+1}- \bm{B}^{k}\| \\
%        &+\, L_{F}\|\bx^{k+2}- \bx^{k+1}\|,
% \end{align*}
%     where the second inequality uses the smoothness of $F(\bx, \bm{B})$ in \Cref{Ass: main assumptions}. Therefore,
%     using the fact $(a + b + c)^{2}\leq 3(a^{2}+ b^{2}+c^{2})$, one has
%     \begin{align}
%          &\|\bm{\Lambda}^{k+1}- \bm{\Lambda}^{k}\|^{2}  \notag \\
%          \leq&\,3\tau^{-2}\|\bm{B}^{k+2}- \bm{B}^{k+1}\|^{2}+ 3(\tau^{-1}+ L_{F})^{2}\|\bm{B}^{k+1}- \bm{B}^{k}\|^{2} \notag \\
%          &+ 3L_{F}^{2}\|\bx^{k+2}- \bx^{k+1}\|^{2},\label{Eq: dual boundedness in proof}
%     \end{align}and completes the proof. 
% \end{proof}
% Next, we analyze the evolution of the Lagrangian function \(\calL\) during the updates of both primal and dual variables. Although the updates of \(\bB\) and \(\bLmbd\) are shared between Algorithms \ref{alg:our_algorithm_sgdm} and \ref{alg:our_algorithm_storm}, the stochastic updates of \(\bx^k\) introduce additional analytical challenges. As indicated by \eqref{Eq: dual boundedness}, the coupling between variables implies that randomness in the \(\bx\)-updates propagates to the updates of \(\bB\) and \(\bLmbd\). To address these challenges, we carefully design the parameter settings for different stochastic estimators to control the variance throughout the optimization process. This leads us to construct the following algorithm-specific Lyapunov functions:
% \[
% \Psi_{\text{sgdm}}^k = \calL(\bx^{k}, \bm{B}^{k}, \bLmbd^{k}) - C_{1}\Delta_{\bB}^{k+1} + C_{2}\Delta_{\bB}^{k} - C_{3}\Delta_{\bm{x}}^{k+1} + C_{4}\Delta_{\bm{x}}^{k}
% \]
% and
% \[
% \Psi_{\text{storm}}^k = \calL(\bx^{k}, \bm{B}^{k}, \bLmbd^{k}) - C_{1}\Delta_{\bB}^{k+1} + C_{2}\Delta_{\bB}^{k} - C_{3}\Delta_{\bm{x}}^{k+1},
% \]
% where 
% \[
% \Delta_{\bx}^{k} = \frac{\|\bx^{k} - \bx^{k-1}\|^{2}}{2}, \quad \Delta_{\bm{B}}^{k} = \frac{\|\bm{B}^{k} - \bm{B}^{k-1}\|^{2}_{F}}{2},
% \]
% and \(C_1, \dots, C_4\) are positive constants specified in the appendix.
% % \(\Psi_{sgdm}\) for \Cref{alg:our_algorithm_sgdm} and \(\Psi_{storm}\) for \Cref{alg:our_algorithm_storm}, and subsequently establish each approximate descent in expectation.
% % % \red{Put the Lyapunov functions here? may be somehow long?}
% % % a Lyapunov sequence specific to 
% % \begin{align}
% %      & \Psi_{storm}(\mathcal{Q}^{k}) = \calL(\bx^{k}, \bm{B}^{k}, \bLmbd^{k}) - C_{1}\Delta_{\bB}^{k+1}+ C_{2}\Delta_{\bB}^{k}- C_{3}\Delta_{\bm{x}}^{k+1}, \\
% %      & \{ \calQ^{k}\}_{k \in \NN}= \{\bx^{k}, \bB^{k}, \bm{\Lambda}^{k}, \bB^{k+1}, \bB^{k-1}, \bx^{k+1}\}_{k \in \NN}\label{Eq: the Lyapunov sequence in STORM},
% % \end{align} 
% % \Cref{alg:our_algorithm_storm}:

% \subsection{Oracle Complexity Results}
% We now present the complexity results of the two algorithms.\footnote{Due to space constraints, necessary lemmas with detailed proofs, as well as proofs of theorems and parameter derivations, are provided in Supplementary Material B.}

% \paragraph{Oracle complexity of \Cref{alg:our_algorithm_sgdm}:}
% We assume the parameters \(\{\eta_k \}\), \(\{\alpha_k, \beta_k\}\), \( \tau \) and \( b_k\) in \Cref{alg:our_algorithm_sgdm} 
% are set as follows:
% \begin{align}
% \label{Eq: the parameters conditions of dualhash-sgdm}
% \eta_k = \frac{\eta}{L_F},\quad \alpha_k = \alpha, \quad \beta_k = \beta,\quad \tau \leq \frac{\sqrt{\tilde{\delta}}}{L_F}, \quad b_k = b = c_bT ,
% \end{align}where \(\eta,\, \alpha, \, \beta,\, \tau,\,\tilde{\delta},\, c_b\) are given constants independent of \(T\) with some \(\nu\), \(\delta\), \( c_{\delta} > 0\).
% % (specifically, 
% % \(\eta =\frac{(1 - 2\alpha - \nu)/2}{\frac{\bar{L}_F}{L_F} + \frac{(1 + \nu)\beta^2}{(1 - \alpha)\alpha}} \) , \(\bar{L}_F = 2L_F + \frac{\tau L_F^2}{\delta} + 3 \delta \tau L_F^2 \),  \(\alpha \in ( 0, \frac{1 - \nu}{2})\), \(\beta \in (0,1)\), \(\tilde{\delta}= c_{\delta} \delta^2\) )
% % \vspace{-6mm}
% \begin{theorem}
% \label{Theorem: the oracle complexity in SGDM} 
% Under \Cref{Ass: main assumptions}, let \(\{\bx^k, \bB^k, \bLmbd^k\}_{k \in \NN}\) be the iterate sequence from \Cref{alg:our_algorithm_sgdm} with the parameters satisfying \eqref{Eq: the parameters conditions of dualhash-sgdm}. Then there exists $R$ uniformly selected from $\{2,\dots,T+1\}$ such that \( (\bx^R, \bB^R, \bLmbd^R )\) satisfies \eqref{equ:esp} provided that \(T\) satisfies
% \begin{align}
%     \label{Eq: the iterate number upper bound of algorithm SGDM}
%         T = \left\lceil \frac{L_{\Delta} \Delta_1+ C_{\sigma}\frac{\sigma^2}{c_b}}{\varepsilon^2} \right\rceil,
% \end{align}
% where  \(\Delta_1 = \Psi_{sgdm}^1 - \Psi_{sgdm}^* \) is the initial Lyapunov function value gap and \(L_{\Delta}\), \(C_{\sigma}\) are the positive constant independent of  \(T\), depending on the parameters \((L_F,  \, \eta, \,\alpha,\, \beta,\, \tau, \, \delta, \, \nu)\). Moreover, the total number of SFO is in the order of \(\calO(\varepsilon^{-4})\).
% \end{theorem}
% % \vspace{5pt}

% \paragraph{Oracle complexity of \Cref{alg:our_algorithm_storm}:} 
% Before proceeding, we introduce an additional standard assumption about the smoothness of \(F\) widely used in variance reduction methods \cite{tran2022hybrid,arjevani2023lower,xu2023momentum}.
% \begin{assumption}
% \label{Ass: the mean squared smoothness of F}
%     There exists $L_{F} > 0$ such that for any index $i$ selected uniformly
%     from $\{1,2,\ldots,n\}$, $\bx, \bx' \in \calX$, and $\bB, \bB' \in \calB$, it holds that
% \begin{align}
%     &\mathbb{E}_{i}\left[ \| \nabla F_{i}(\bx,\bB) - \nabla F_{i}(\bx',\bB')\|^{2}\right] \notag \\
%     \leq&\, L^{2}_F \left[ \|\bx - \bx'\|^{2} + \|\bB - \bB'\|^{2} \right].
% \end{align}
% \end{assumption} 

% This assumption is slightly stronger than the \(L_F\)-smoothness of the finite-sum function \(F\) in \Cref{Ass: main assumptions} (ii). In particular, this assumption implies \Cref{Ass: main assumptions} (ii) by Jensen's inequality, but not conversely.

% As shown in \cite{xu2023momentum}, employing a moderately large initial batch size in VR methods improves the order of complexity. Motivated by this, we sample \(b_1 = \mathcal{O}(T^{1/3})\) times at the initial point to achieve reduced complexity:
% \begin{align}
% \label{Eq: the batch size selection of VR estimator}
% b_k=
% \begin{cases}
% b_1,  &\text{if }k =1,\\
% b, &\text{if } k > 1.
% \end{cases}
% \end{align}
% To ensure the optimal convergence for DualHash-StoRM, we adopt the following parameter setting: 
% % We then assume the parameters \(\{\eta_k \}\), \(\{ \rho_k\}\), \( \tau \) and \(b_k \) are set as follows:
% \begin{align}
% \label{Eq: the parameters conditions in dualhash-storm}
% \eta_k = \frac{\eta}{\tilde{L}_{F}T^{\frac{1}{3}}},\quad \rho_k = \frac{8\rho\eta^2}{T^{\frac{2}{3}}},  \quad \tau \leq \frac{\sqrt{\tilde{\delta}}}{L_F},\quad  b_1 = c_b T^{1/3},
% \end{align} where \(\eta,\, \rho, \tilde{\delta},c_b, b\) are given constants independent of \(T\) with some  \( \delta\), \( c_{\delta} > 0\).
% % (specifically,  \(0< \eta \leq \frac{1}{2}\),  \(\tilde{L}_F = L_F + \frac{\tau L_F^2}{\delta} + 3 \delta \tau L_F^2 \) , \(1 \leq \rho \leq \frac{1}{8\eta^2}\), \(\tilde{\delta} = c_{\delta} \delta^2\) with some  \( \delta\), \( c_{\delta} > 0\).
% \begin{theorem}
% \label{Theorem: the oracle complexity in STORM}
% Under \Cref{Ass: main assumptions} (i), (iii) and \Cref{Ass: the mean squared smoothness of F}, let \(\{\bx^k, \bB^k, \bLmbd^k\}_{k \in \NN}\) be the iterate sequence from \Cref{alg:our_algorithm_storm} with the parameters satisfying \eqref{Eq: the parameters conditions in dualhash-storm}. Then there exists $R$ uniformly selected from $\{2,\dots,T+1\}$ such that \( (\bx^R, \bB^R, \bLmbd^R )\) satisfies \eqref{equ:esp} provided that  \(T\)  satisfies
% \begin{align}
%     \label{Eq: the iterate number upper bound of algorithm VR}
%     % T = \left\lceil\frac{\left[\left(\frac{2(3\tilde{\delta} + 5) \tilde{L}_{F}}{\eta} + ( 7 + 3\tilde{\delta}) \right) \left(\Psi(\calQ^1) - \Psi^*\right) + \frac{7 + 3\tilde{\delta}\sigma^2}{2\rho \eta^2c_b} + \frac{32(7 + 3\tilde{\delta})\rho\eta^2\sigma^2}{b} \right]^{\frac{3}{2}}}{\varepsilon^3} \right\rceil.
%         T = \left\lceil \frac{\left(\tilde{L}_{\Delta} \tilde{\Delta}_1+ \tilde{C}_{\sigma}\left( \frac{\sigma^2}{ 4\rho \eta^2 c_b} + \frac{32\rho\eta^2\sigma^2}{b} \right)\right)^{\frac{3}{2}}}{\varepsilon^3} \right\rceil,
% \end{align}where \(\tilde{\Delta}_1 = \Psi_{storm}^1 - \Psi_{storm}^* \) is the initial Lyapunov function value gap and \(\tilde{L}_{\Delta}\), \(\tilde{C}_{\sigma}\) are the positive constant independent of \(T\), depending on the parameters \((L_F, \, \eta, \, \rho, \tau, \, \delta)\). Moreover, the total number of SFO is in the order of \(\calO(\varepsilon^{-3})\). 
% \end{theorem} 

% \begin{remark}
% % \begin{itemize}[label=(\roman*), leftmargin=2em]
% \begin{enumerate}[label=(\roman*), leftmargin=2em]
% \item The results in Theorems \ref{Theorem: the oracle complexity in SGDM} and \ref{Theorem: the oracle complexity in STORM} are novel for stochastic, non-convex, multi-block primal-dual optimization, a challenging area where existing work has largely focused on one-block \cite{xu2023momentum,shi2025momentum,tran2022hybrid} or deterministic settings \cite{ADMMNolinear2024}.  
%     \item  \Cref{Theorem: the oracle complexity in SGDM} with constant parameters from \eqref{Eq: the parameters conditions of dualhash-sgdm} attains both oracle and sample \(\mathcal{O}(\varepsilon^{-4})\) complexities, which matches that of stochastic first-order methods (SFOMs) for one-block stochastic optimization \cite{Ghadimi2013MinibatchSA}. Meanwhile, \Cref{Theorem: the oracle complexity in STORM} with the refined constant parameters from \eqref{Eq: the parameters conditions in dualhash-storm} attains both oracle and sample $\mathcal{O}(\varepsilon^{-3})$ complexities. This matches the best-known lower bound \cite{arjevani2023lower} and the complexity of variance-reduced SFOMs for single-block problems \cite{shi2025momentum,xu2023momentum}. Notably, despite tackling a multi-block problem, our primal-dual framework preserves this efficiency without degrading the dependence on $\varepsilon$.
%    \item  While DualHash-StoRM (\Cref{alg:our_algorithm_storm}) with parameters in \eqref{Eq: the parameters conditions in dualhash-storm} achieves better complexity than DualHash-StoM (\Cref{alg:our_algorithm_sgdm}) under \eqref{Eq: the parameters conditions of dualhash-sgdm}, the latter relies on a weaker assumption (\Cref{Ass: main assumptions} [ii]). In our experiments, while DualHash-StoRM enjoys faster training convergence, DualHash-StoM enerally delivers better final retrieval performance.
% \end{enumerate}
% \end{remark}
\section{Complexity analysis}
\label{Section: Convergence analysis}

We now investigate the complexity analysis for DualHash-StoM (\Cref{alg:our_algorithm_sgdm}) and DualHash-StoRM (\Cref{alg:our_algorithm_storm}). Our objective is to find an approximate stochastic critical point of the Lagrangian function \eqref{Eq: simplified_lagrangian}, i.e., a point satisfying 
\begin{equation}\label{equ:esp}
\mathbb{E} [\text{dist}^2(\bm{0}, \partial \calL (\bx^*, \bm{B}^*, \bm{\Lambda}^* ))] \leq \varepsilon^2
\end{equation}
for a given \(\varepsilon > 0\). We assume in this paper the set \(
\text{crit} \calL := \{(\bx^*, \bm{B}^*, \bm{\Lambda}^* ) \in \mathcal{X} \times \mathcal{B} \times \mathcal{Y^*}: \bm{0} \in \partial \mathcal{L}(\bx^*, \bm{B}^*, \bm{\Lambda}^*)\}
\) is nonempty. This focus is standard in non-convex stochastic optimization analysis, as finding a global or even local minimizer for non-convex optimization problems is generally NP-hard. In particular, we can prove that a critical point of the Lagrangian function \eqref{Eq: simplified_lagrangian} is a first-order necessary optimal condition of problem \eqref{Prop: block-wise regularized problem} as follows. 

\begin{lemma}
\label{Lemma: Optimality condition}
    Suppose that \(\min\limits_{\bx, \bB} \calL(\bx, \bB, \bLmbd) > -\infty\)  and \((\bx^*, \bB^*, \bLmbd^*)\) are the optimal solutions of the problem \eqref{Prop: block-wise regularized problem} and the problem \eqref{Prop: the dual problem}, respectively. Then, it holds that \((\bx^*, \bB^*, \bLmbd^*) \in \text{crit} \calL\), i.e., \(\bm{0} \in \partial \mathcal{L}(\bx^*, \bB^*, \bLmbd^*)\).
\end{lemma} 
\setcounter{proof}{0}
\begin{proof}
    According to \cite[Theorem 2.158]{temporary-citekey-11039},
    $(\bx^{*}, \bm{B}^{*}, \bm{\Lambda}^{*})$ are optimal solutions to the
    primal problem \eqref{Prop: block-wise regularized problem} and the dual problem
    \eqref{Prop: the dual problem} if and only if:
    \begin{align}
        \label{Eq: the sufficient and necessary condition}
        \begin{cases}
        (\bx^{*}, \bm{B}^{*}) \in \argmin\limits_{\bx \in \calX, \bm{B} \in \mathcal{B}}\mathcal{L}(\bx^{*}, \bm{B}^{*}, \bm{\Lambda}^{*}), \\ 
        0 = h(\bm{B}^{*}) + h^{*}(\bm{\Lambda}^{*}) - \langle \bm{\Lambda}^{*}, \bm{B}^{*}\rangle.
        \end{cases}
    \end{align}

    The first condition in \eqref{Eq: the sufficient and necessary condition} implies
    that $\bm{0}$ belongs to the subdifferential of $\mathcal{L}$ with respect to
    $(\bx, \bm{B})$ at $(\bx^{*}, \bm{B}^{*}, \bm{\Lambda}^{*})$. Since $\calL$
    is continuously differentiable in $(\bx, \bm{B})$, we have
    \begin{align}
        \label{Eq: the sufficient and necessary condition for primal}
        \begin{cases}
        \bm{0} = \nabla f(\bx^{*}) + \nabla_{\bx}G(\bx^{*}, \bm{B}^{*}),\\ 
        \bm{0} = \nabla_{\bB}G(\bx^{*}, \bm{B}^{*}) + \bm{\Lambda}^{*}.
        \end{cases}
    \end{align}
    For the second condition in \eqref{Eq: the sufficient and necessary condition},
    by the definition of the conjugate function, if this equality holds, then $\bm{\Lambda}^{*}$ is the optimal solution of $\sup \{ \left<\bb^{*}, \blmbd\right> - h^{*}(\blmbd) \}$, i.e.,
    \begin{align}
        \label{Eq: the sufficient and necessary condition for dual}\bm{B}^{*}\in \partial h^{*}(\bm{\Lambda}^{*}).
    \end{align}
    Combining \eqref{Eq: the sufficient and necessary condition for primal} and \eqref{Eq: the sufficient and necessary condition for dual}, we obtain $\bm{0}\in \partial \mathcal{L}(\bx^{*}, \bm{B}^{*}, \bm{\bLmbd}^{*})$, completing the proof. 
    % 
\end{proof}

We then make the following assumptions generally used in stochastic optimization \cite{Ghadimi2013MinibatchSA,huangMiniBatchStochasticADMMs2019,xu2023momentum}.
\begin{assumption} 
\label{Ass: main assumptions}
\begin{itemize}[leftmargin=2em]
    \item[{(i)}] We assume that the sequences \( \{(\bx^k, \bm{B}^k, \bm{\Lambda}^k) \}_{k \in \NN} \) generated by Algorithms \ref{alg:our_algorithm_sgdm} and \ref{alg:our_algorithm_storm} are bounded. Moreover, the objective of  \eqref{Prop: block-wise regularized problem} is lower bounded by \(J^*\).
    \item[{(ii)}] We assume that \(F(\bx, \bB)\) is $L_F$-smooth with respect to $(\bx, \bB)$, where $L_F = \mathcal{O}\left(1\right)$.
    \item[{(iii)}] Suppose index $i \sim \text{Uniform}\{1, \ldots, n\}$.  For any \(\bB \in \calB\), there exists a constant $\sigma^2 \geq 0$ such that $\mathbb{E}_{i}[\nabla_{\bx} F_i(\bx, \bb_i)] = \nabla_{\bx} F(\bx, \bB)$ and $\mathbb{E}_{i}[\|\nabla_{\bx} F_i(\bx, \bb_i) - \nabla_{\bx} F(\bx, \bB)\|^2] \leq \sigma^2$. 
\end{itemize}
\end{assumption}
\begin{remark}
\begin{enumerate}[label=(\roman*), leftmargin=2em]
    \item The boundedness in \Cref{Ass: main assumptions} (i) is maintained in practice for each variable:  $\bm{\Lambda}^k$ remains bounded due to its proximal update rule \eqref{Eq: the update of lambda} for $\bm{\Lambda}^k$ with the truncation property of $\text{prox}_{h^*}$ in \eqref{Eq: the proximal mapping of conjugate function}; The $\bB^k$ naturally converges toward $\pm 1$ through W-type regularization; and network parameters $\bx^k$ are stabilized via weight decay and pretrained initialization, ensuring objective coercivity.
    \item The smoothness constant \(L_F\) remains \(\mathcal{O}(1)\), as the penalty formulation is designed to cancel the potential \(\mathcal{O}(\sqrt{n})\) scaling of individual components. A rigorous analysis is provided in Appendix C.
\end{enumerate}
\end{remark}
Our analysis aims to bound \( \EE [\| \partial \calL\|^2]\). However, due to the existence of dual variables \(\bLmbd\), the Lagrangian function \eqref{Prop: the dual problem} lacks the descent property. This makes it difficult to apply standard approaches from prior work \cite{tran2022hybrid,wang2019spiderboost,xu2023momentum,zhang2022stochastic}, which often leverage the descent of their objective functions or corresponding potential functions. Instead, we shall construct an appropriate Lyapunov function. To enable this construction, we first analyze how to control the dual variable dynamics by establishing bounds on \(\|\bm{\Lambda}^{k+1} - \bm{\Lambda}^k \|^2\) with the primal variables residuals.
\begin{lemma}
\label{Lemma: dual boundedness}
Under \Cref{Ass: main assumptions}, for \(k \geq 1 \), it holds that
    \begin{align}
        &\|\bm{\Lambda}^{k+1}- \bm{\Lambda}^k \|^2 \notag \\
        \leq&\,3 \tau^{-2} \|\bB^{k+2} - \bB^{k+1}\|^2 + 3(\tau^{-1} + L_F)^2  \|\bB^{k+1} - \bB^{k}\|^2 + 3L_F^2  \|\bx^{k+2} - \bx^{k+1}\|^2.
         \label{Eq: dual boundedness}
    \end{align}
\end{lemma}
\begin{proof}
    Using the update of $\bm{B}^{k}$ in \eqref{Eq: the update of B} twice yields that
\begin{align*}
    &\|\bm{\Lambda}^{k+1}- \bm{\Lambda}^{k}\| \notag \\
    =&\, \left\| \tau^{-1}(\bm{B}^{k+1}- \bm{B}^{k+2}) + \tau^{-1}(\bm{B}^{k+1}- \bm{B}^{k}) + \nabla_{\bB}F(\bx^{k+1}, \bm{B}^{k}) - \nabla_{\bB}F(\bx^{k+2}, \bm{B}^{k+1}) \right\| \\
    \leq&\, \tau^{-1}\|(\bm{B}^{k+1}- \bm{B}^{k+2})\| + (\tau^{-1}+ L_{F}) \|\bm{B}^{k+1}- \bm{B}^{k}\| + L_{F}\|\bx^{k+2}- \bx^{k+1}\|,
\end{align*}
    where the second inequality uses the smoothness of $F(\bx, \bm{B})$ in \Cref{Ass: main assumptions}. Therefore, using the fact $(a + b + c)^{2}\leq 3(a^{2}+ b^{2}+c^{2})$, one has
    \begin{align}
         &\|\bm{\Lambda}^{k+1}- \bm{\Lambda}^{k}\|^{2} \notag \\
         \leq&\, 3\tau^{-2}\|\bm{B}^{k+2}- \bm{B}^{k+1}\|^{2}+ 3(\tau^{-1}+ L_{F})^{2}\|\bm{B}^{k+1}- \bm{B}^{k}\|^{2} + 3L_{F}^{2}\|\bx^{k+2}- \bx^{k+1}\|^{2},\label{Eq: dual boundedness in proof}
    \end{align}
    and completes the proof. 
\end{proof}
Next, we analyze the evolution of the Lagrangian function \(\calL\) during the updates of both primal and dual variables. Although the updates of \(\bB\) and \(\bLmbd\) are shared between Algorithms \ref{alg:our_algorithm_sgdm} and \ref{alg:our_algorithm_storm}, the stochastic updates of \(\bx^k\) introduce additional analytical challenges. As indicated by \eqref{Eq: dual boundedness}, the coupling between variables implies that randomness in the \(\bx\)-updates propagates to the updates of \(\bB\) and \(\bLmbd\). To address these challenges, we carefully design the parameter settings for different stochastic estimators to control the variance throughout the optimization process. This leads us to construct the following algorithm-specific Lyapunov functions:
\[
\Psi_{\text{sgdm}}^k = \calL(\bx^{k}, \bm{B}^{k}, \bLmbd^{k}) - C_{1}\Delta_{\bB}^{k+1} + C_{2}\Delta_{\bB}^{k} - C_{3}\Delta_{\bm{x}}^{k+1} + C_{4}\Delta_{\bm{x}}^{k}
\]
and
\[
\Psi_{\text{storm}}^k = \calL(\bx^{k}, \bm{B}^{k}, \bLmbd^{k}) - C_{1}\Delta_{\bB}^{k+1} + C_{2}\Delta_{\bB}^{k} - C_{3}\Delta_{\bm{x}}^{k+1},
\]
where 
\[
\Delta_{\bx}^{k} = \frac{\|\bx^{k} - \bx^{k-1}\|^{2}}{2}, \quad \Delta_{\bm{B}}^{k} = \frac{\|\bm{B}^{k} - \bm{B}^{k-1}\|^{2}_{F}}{2},
\]
and \(C_1, \dots, C_4\) are positive constants specified in the appendix.

\subsection{Oracle Complexity Results}
We now present the complexity results of the two algorithms.\footnote{Necessary lemmas with detailed proofs, as well as proofs of theorems and parameter derivations, are provided in Appendix.}

\paragraph{Oracle complexity of \Cref{alg:our_algorithm_sgdm}:}
We assume the parameters \(\{\eta_k \}\), \(\{\alpha_k, \beta_k\}\), \( \tau \) and \( b_k\) in \Cref{alg:our_algorithm_sgdm} 
are set as follows:
\begin{align}
\label{Eq: the parameters conditions of dualhash-sgdm}
\eta_k = \frac{\eta}{L_F},\quad \alpha_k = \alpha, \quad \beta_k = \beta,\quad \tau \leq \frac{\sqrt{\tilde{\delta}}}{L_F}, \quad b_k = b = c_bT ,
\end{align}
where \(\eta,\, \alpha, \, \beta,\, \tau,\,\tilde{\delta},\, c_b\) are given constants independent of \(T\) with some \(\nu\), \(\delta\), \( c_{\delta} > 0\).
\begin{theorem}
\label{Theorem: the oracle complexity in SGDM} 
Under \Cref{Ass: main assumptions}, let \(\{\bx^k, \bB^k, \bLmbd^k\}_{k \in \NN}\) be the iterate sequence from \Cref{alg:our_algorithm_sgdm} with the parameters satisfying \eqref{Eq: the parameters conditions of dualhash-sgdm}. Then there exists $R$ uniformly selected from $\{2,\dots,T+1\}$ such that \( (\bx^R, \bB^R, \bLmbd^R )\) satisfies \eqref{equ:esp} provided that \(T\) satisfies
\begin{align}
    \label{Eq: the iterate number upper bound of algorithm SGDM}
        T = \left\lceil \frac{L_{\Delta} \Delta_1+ C_{\sigma}\frac{\sigma^2}{c_b}}{\varepsilon^2} \right\rceil,
\end{align}
where  \(\Delta_1 = \Psi_{sgdm}^1 - \Psi_{sgdm}^* \) is the initial Lyapunov function value gap and \(L_{\Delta}\), \(C_{\sigma}\) are the positive constant independent of  \(T\), depending on the parameters \((L_F,  \, \eta, \,\alpha,\, \beta,\, \tau, \, \delta, \, \nu)\). Moreover, the total number of SFO is in the order of \(\calO(\varepsilon^{-4})\).
\end{theorem}

\paragraph{Oracle complexity of \Cref{alg:our_algorithm_storm}:} 
Before proceeding, we introduce an additional standard assumption about the smoothness of \(F\) widely used in variance reduction methods \cite{tran2022hybrid,arjevani2023lower,xu2023momentum}.
\begin{assumption}
\label{Ass: the mean squared smoothness of F}
    There exists $L_{F} > 0$ such that for any index $i$ selected uniformly from $\{1,2,\ldots,n\}$, $\bx, \bx' \in \calX$, and $\bB, \bB' \in \calB$, it holds that
\begin{align}
    \mathbb{E}_{i}\left[ \| \nabla F_{i}(\bx,\bB) - \nabla F_{i}(\bx',\bB')\|^{2}\right] \leq L^{2}_F \left[ \|\bx - \bx'\|^{2} + \|\bB - \bB'\|^{2} \right].
\end{align}
\end{assumption} 

This assumption is slightly stronger than the \(L_F\)-smoothness of the finite-sum function \(F\) in \Cref{Ass: main assumptions} (ii). In particular, this assumption implies \Cref{Ass: main assumptions} (ii) by Jensen's inequality, but not conversely.

As shown in \cite{xu2023momentum}, employing a moderately large initial batch size in VR methods improves the order of complexity. Motivated by this, we sample \(b_1 = \mathcal{O}(T^{1/3})\) times at the initial point to achieve reduced complexity:
\begin{align}
\label{Eq: the batch size selection of VR estimator}
b_k=
\begin{cases}
b_1,  &\text{if }k =1,\\
b, &\text{if } k > 1.
\end{cases}
\end{align}
To ensure the optimal convergence for DualHash-StoRM, we adopt the following parameter setting: 
\begin{align}
\label{Eq: the parameters conditions in dualhash-storm}
\eta_k = \frac{\eta}{\tilde{L}_{F}T^{\frac{1}{3}}},\quad \rho_k = \frac{8\rho\eta^2}{T^{\frac{2}{3}}},  \quad \tau \leq \frac{\sqrt{\tilde{\delta}}}{L_F},\quad  b_1 = c_b T^{1/3},
\end{align} 
where \(\eta,\, \rho, \tilde{\delta},c_b, b\) are given constants independent of \(T\) with some  \( \delta\), \( c_{\delta} > 0\).
\begin{theorem}
\label{Theorem: the oracle complexity in STORM}
Under \Cref{Ass: main assumptions} (i), (iii) and \Cref{Ass: the mean squared smoothness of F}, let \(\{\bx^k, \bB^k, \bLmbd^k\}_{k \in \NN}\) be the iterate sequence from \Cref{alg:our_algorithm_storm} with the parameters satisfying \eqref{Eq: the parameters conditions in dualhash-storm}. Then there exists $R$ uniformly selected from $\{2,\dots,T+1\}$ such that \( (\bx^R, \bB^R, \bLmbd^R )\) satisfies \eqref{equ:esp} provided that  \(T\)  satisfies
\begin{align}
    \label{Eq: the iterate number upper bound of algorithm VR}
        T = \left\lceil \frac{\left(\tilde{L}_{\Delta} \tilde{\Delta}_1+ \tilde{C}_{\sigma}\left( \frac{\sigma^2}{ 4\rho \eta^2 c_b} + \frac{32\rho\eta^2\sigma^2}{b} \right)\right)^{\frac{3}{2}}}{\varepsilon^3} \right\rceil,
\end{align}
where \(\tilde{\Delta}_1 = \Psi_{storm}^1 - \Psi_{storm}^* \) is the initial Lyapunov function value gap and \(\tilde{L}_{\Delta}\), \(\tilde{C}_{\sigma}\) are the positive constant independent of \(T\), depending on the parameters \((L_F, \, \eta, \, \rho, \tau, \, \delta)\). Moreover, the total number of SFO is in the order of \(\calO(\varepsilon^{-3})\). 
\end{theorem} 

\begin{remark}
\begin{enumerate}[label=(\roman*), leftmargin=2em]
\item The results in Theorems \ref{Theorem: the oracle complexity in SGDM} and \ref{Theorem: the oracle complexity in STORM} are novel for stochastic, non-convex, multi-block primal-dual optimization, a challenging area where existing work has largely focused on one-block \cite{xu2023momentum,shi2025momentum,tran2022hybrid} or deterministic settings \cite{ADMMNolinear2024}.  
    \item  \Cref{Theorem: the oracle complexity in SGDM} with constant parameters from \eqref{Eq: the parameters conditions of dualhash-sgdm} attains both oracle and sample \(\mathcal{O}(\varepsilon^{-4})\) complexities, which matches that of stochastic first-order methods (SFOMs) for one-block stochastic optimization \cite{Ghadimi2013MinibatchSA}. Meanwhile, \Cref{Theorem: the oracle complexity in STORM} with the refined constant parameters from \eqref{Eq: the parameters conditions in dualhash-storm} attains both oracle and sample $\mathcal{O}(\varepsilon^{-3})$ complexities. This matches the best-known lower bound \cite{arjevani2023lower} and the complexity of variance-reduced SFOMs for single-block problems \cite{shi2025momentum,xu2023momentum}. Notably, despite tackling a multi-block problem, our primal-dual framework preserves this efficiency without degrading the dependence on $\varepsilon$.
   \item  While DualHash-StoRM (\Cref{alg:our_algorithm_storm}) with parameters in \eqref{Eq: the parameters conditions in dualhash-storm} achieves better complexity than DualHash-StoM (\Cref{alg:our_algorithm_sgdm}) under \eqref{Eq: the parameters conditions of dualhash-sgdm}, the latter relies on a weaker assumption (\Cref{Ass: main assumptions} [ii]). In our experiments, while DualHash-StoRM enjoys faster training convergence, DualHash-StoM enerally delivers better final retrieval performance.
\end{enumerate}
\end{remark}